\chapter*{一、緒論}
\section*{1.1 研究背景與動機}

時間序列是種常見的資料型態，過去觀察值通常與未來變化具有關聯，如何利用既有資料預測尚未發生的數值，長期以來一直是統計學與資料分析的重要研究領域。時間序列預測不僅希望掌握資料未來的變動方向，也可作為決策與風險管理的依據。

過去常見的時間序列預測方法以統計模型為主。例如，ARIMA模型利用序列本身的落後值與歷史誤差描述線性時間相依關係，是時間序列分析中廣泛使用的方法\cite{boxjenkins2015}。當資料包含無法由單一線性模型充分描述的非線性變化時，人工神經網路可藉由隱藏層與非線性激活函數學習較複雜的輸入與輸出關係。另有研究將ARIMA與ANN結合，由ARIMA處理線性結構，再由ANN補充剩餘的非線性訊號，以提高不同型態時間序列的預測能力\cite{zhang2003,khashei2011}。

\cite{lin2024thesis}進一步使用Kalman filter將時間序列分解為變動較平緩的低頻成分與變動較快速的高頻成分，分別以ARIMA與ANN預測，並利用Sort Bootstrap建立預測區間。然而，只提供單一點預測無法表達預測結果的不確定性。即使點預測接近真實值，使用者仍無法得知可能的誤差範圍；反之，預測區間若為了提高覆蓋率而開得過寬，也會降低實際參考價值。因此，區間預測除了考慮真實值是否落在區間內，也必須同時衡量區間寬度及未覆蓋時的偏離程度。\cite{hung2026thesis}沿用\cite{lin2024thesis}分解架構，使用更進一步的方法來形成點預測及區間預測。這些方法顯示，面對同一條時間序列，可先分離不同變動特性的成分，再選擇適合的模型進行預測。

但是時間序列的預測區間還面臨兩項問題。第一，預測某一期時只能使用該期以前的資訊；若用來決定區間的殘差不是依時間順序產生，可能無法代表真正的樣本外預測誤差。第二，預測誤差的尺度可能隨時間改變。當序列位於平穩、低波動狀態時，所需區間可能較窄；當序列快速變動或進入高波動狀態時，所需區間則可能變寬。若所有時間點都使用同一組未經尺度調整的殘差，便可能在部分狀態下區間過寬，在另一些狀態下又不足以涵蓋真實值。

為處理上述問題，本研究提出Out-of-Time Adaptive EnbPI（以下簡稱OOT-AEnbPI）。此方法在訓練期間依照時間順序建立多個future blocks，每次只使用較早資料訓練，再預測較晚的資料，以取得較接近實際樣本外情境的殘差。接著，利用訓練期殘差表現選擇適當的residual window，並從殘差經驗分布建立最短的非對稱預測區間。

在OOT-AEnbPI的核心方法之外，本研究進一步加入Local-scale調整，形成OOT-AEnbPI-LS。此版本根據當期預測水準、變化方向及短期與長期波動，尋找歷史上狀態相近的時間點，以估計當期局部誤差尺度，使預測區間可隨資料狀態調整。

本研究將OOT-AEnbPI與\cite{hung2026thesis}及\cite{lin2024thesis}所提出的方法進行比較，並同時檢驗各方法的原始區間與Local-scale區間。此外，本文保留OOT-AEnbPI的點預測，將IFOMC狀態轉移概念應用於其OOT residual，形成OOT-AEnbPI-IFOMC，作為另一種區間預測器進行比較。實驗包含GARCH模擬資料、0050價格資料、太陽黑子月資料及風速資料。結果顯示，Local-scale整體上可使多數方法的coverage更接近名目水準，並改善interval score。在相同OOT-AEnbPI點預測下，OOT-AEnbPI-LS於四組資料的三種區間比較中皆取得最低interval score，顯示依時間順序建立殘差並依局部狀態調整尺度，有助於提升預測區間的實用性。OOT-AEnbPI-IFOMC在GARCH及風速資料的部分指標上小幅改善原始區間，但未全面優於OOT-AEnbPI-LS。

\section*{1.2 論文架構}

後續章節安排如下：第二章整理各項模型及區間方法的理論背景；第三章說明本文提出之OOT-AEnbPI，包括Kalman高低頻分解、ARIMA--ANN ensemble點預測器、OOT residual、自適應residual window及最短非對稱區間，並介紹OOT-AEnbPI-LS與OOT-AEnbPI-IFOMC兩種區間版本，其後說明五種比較方法及評估指標；第四章呈現四組資料的點預測與區間預測結果；最後說明研究結論、限制與未來研究方向。
